% =====================================================================
\section{Ответы на вопросы «Проверь себя»}
% =====================================================================
Сначала ответьте сами и только потом сверяйтесь. Ответы краткие: полное объяснение есть в тексте урока.

\begingroup
\newcommand{\anslesson}[1]{\par\medskip\noindent\textsf{\bfseries\color{accent}#1}\par\nopagebreak}
\setlist[enumerate,1]{itemsep=1pt,topsep=2pt,leftmargin=*}
\small

\anslesson{Вводная глава «Сигналы»}
\begin{enumerate}
\item Выключатель передаёт два состояния; ручка громкости задаёт плавно меняющуюся величину.
\item 1,5 В находится в неопределённой области: гарантированного цифрового уровня нет.
\item $2^5=32$ варианта.
\item $32+8+4+1=45$.
\end{enumerate}
\anslesson{Вводная глава «Выводы платы»}
\begin{enumerate}
\item Это короткое замыкание: ток ограничен почти только сопротивлением проводов и источника.
\item В опытах выбираем ADC1, GPIO1--GPIO10; ADC2 используется также радиомодулем.
\item Мотор требует большего тока, чем может дать GPIO; нужен драйвер или транзистор с защитой.
\end{enumerate}

\anslesson{Урок 0. Знакомство с ESP32-S3}
\begin{enumerate}
  \item Чип --- сама микросхема; модуль --- чип вместе с Flash-памятью, кварцем и антенной. N16R8 означает 16 МБ Flash и 8 МБ PSRAM.
  \item «UART» подключён через микросхему-преобразователь: это самый надёжный путь для прошивки и сообщений. «USB» идёт прямо к чипу и позволяет плате быть USB-устройством и отлаживаться через JTAG.
  \item Войти в загрузчик вручную: зажать BOOT, нажать и отпустить RST, отпустить BOOT, запустить загрузку, после неё нажать RST.
\end{enumerate}

\anslesson{Урок 1. Цифровой выход}
\begin{enumerate}
  \item Резистор ограничивает ток, иначе сгорит светодиод или вывод. $R = (U_\text{вых} - U_F) / I$, например $(3{,}3 - 2{,}0)/0{,}006 \approx 220$ Ом.
  \item Во время \code{delay()} программа стоит на месте и не может читать кнопки, датчики и сеть.
  \item Вычитание беззнаковых чисел даёт верную разницу даже после переполнения счётчика, а сравнение \code{now > last + INTERVAL} --- нет.
\end{enumerate}

\anslesson{Урок 2. Последовательный порт}
\begin{enumerate}
  \item У сторон нет общего тактового сигнала, длительность бита задана заранее выбранной скоростью. При разных скоростях биты читаются не в те моменты, и получается «мусор».
  \item TX одного устройства соединяют с RX другого и наоборот. Общая земля нужна, потому что напряжение сигнала измеряется относительно неё.
  \item Чтобы не ждать данных, которых ещё нет: программа читает только пришедшее и сразу продолжает работу.
\end{enumerate}

\anslesson{Урок 3. Цифровой вход}
\begin{enumerate}
  \item Вход, никуда не подключённый, ловит помехи и читается случайно. Подтягивающий резистор задаёт уровень по умолчанию.
  \item Из-за дребезга: контакты при замыкании несколько раз подпрыгивают, и программа видит серию нажатий.
  \item Можно изменить переменную или установить флаг. Нельзя вызывать \code{delay()}, \code{Serial.print()} и другие медленные функции.
\end{enumerate}

\anslesson{Урок 4. ШИМ}
\begin{enumerate}
  \item Это доля периода, когда сигнал в HIGH. Среднее напряжение равно этой доле, умноженной на \SI{3.3}{V}.
  \item Чем выше частота, тем меньше доступное разрешение: $f_\text{max} \approx \SI{80}{MHz} / 2^\text{бит}$.
  \item Глаз воспринимает яркость нелинейно; без коррекции яркость быстро растёт в начале шкалы и почти не меняется в конце.
\end{enumerate}

\anslesson{Урок 5. АЦП}
\begin{enumerate}
  \item Она учитывает заводскую калибровку чипа и сразу возвращает милливольты.
  \item ADC2 используется также Wi-Fi, и доступ зависит от режима драйвера; для курса выбираем ADC1 (\pin{1}--\pin{10}).
  \item Через делитель из двух резисторов, который понижает напряжение до допустимого. Конденсатор сглаживает помехи на входе АЦП.
\end{enumerate}

\anslesson{Урок 6. Шина I\textsuperscript{2}C}
\begin{enumerate}
  \item Устройства могут только притягивать линию к земле; в HIGH её возвращают подтягивающие резисторы.
  \item Адреса всех ответивших устройств. Если ответа нет --- ошибка в питании, проводах, выводах или адресе.
  \item Отправка буфера занимает десятки миллисекунд; постоянная перерисовка затормозит всю программу.
\end{enumerate}

\anslesson{Урок 7. FreeRTOS}
\begin{enumerate}
  \item У задачи свой бесконечный цикл, стек и приоритет; планировщик сам переключает задачи. \code{loop()} --- это одна из таких задач.
  \item Очередь --- чтобы передавать данные от одной задачи другой. Мьютекс --- чтобы защитить общее состояние, которое читают и меняют несколько задач.
  \item По перезагрузке с сообщением \code{Stack canary} и по малому остатку, который показывает \code{uxTaskGetStackHighWaterMark()}.
\end{enumerate}

\anslesson{Урок 8. Wi-Fi и веб-сервер}
\begin{enumerate}
  \item Станция подключается к существующей сети роутера; точка доступа сама создаёт сеть, к которой подключаются другие.
  \item Браузер отправляет запрос GET /toggle, плата выполняет обработчик и отвечает перенаправлением 303, браузер заново запрашивает главную страницу.
  \item Чтобы программа не зависла навсегда, если сети нет или пароль неверен.
  \item GET читает ресурс, POST отправляет данные или команду. HTTP не шифрует ни один метод; для шифрования нужен HTTPS (урок 13).
  \item HTML задаёт структуру и содержание, CSS --- оформление, JavaScript --- поведение в браузере.
  \item \code{border-radius: 50\%;} при одинаковых ширине и высоте элемента.
\end{enumerate}

\anslesson{Урок 9. Сон}
\begin{enumerate}
  \item После light-sleep программа продолжается со следующей строки; после deep-sleep чип перезапускается и \code{setup()} выполняется заново.
  \item Переменные в RTC-памяти, объявленные с \code{RTC\_DATA\_ATTR}.
  \item От длительности и тока активной фазы (особенно работы Wi-Fi): ток сна на несколько порядков меньше.
\end{enumerate}

\anslesson{Урок 10. SPI}
\begin{enumerate}
  \item У шины нет адресов: ведущий выбирает устройство линией CS, и отвечать вправе только выбранное.
  \item Длинные провода и макетная плата искажают фронты сигналов; на высокой частоте появляются ошибки.
  \item До закрытия часть данных остаётся в буфере и пропадёт при отключении питания.
\end{enumerate}

\anslesson{Урок 11. Память}
\begin{enumerate}
  \item Счётчик пробуждений --- в RTC-памяти, пароль --- в NVS, большой буфер --- в PSRAM.
  \item Ресурс ячеек Flash ограничен; частая запись быстро их изнашивает.
  \item \code{psramFound()} должен вернуть \code{true}, а указатель после \code{ps\_malloc()} не должен быть \code{nullptr}.
\end{enumerate}

\anslesson{Урок 12. ESP-NOW, MQTT, BLE}
\begin{enumerate}
  \item Пульт без роутера --- ESP-NOW; умный дом --- MQTT; настройка с телефона --- BLE.
  \item Брокер хранит последнее такое сообщение и сразу отдаёт его каждому новому подписчику.
  \item Сетевые функции медленные и могут блокироваться; в обработчике это приводит к срабатыванию watchdog и сбоям.
\end{enumerate}

\anslesson{Урок 13. HTTPS и OTA}
\begin{enumerate}
  \item У сертификата есть срок действия; не зная текущей даты, устройство не может его проверить.
  \item Соединение шифруется, но подлинность сервера не проверяется: устройство не отличит настоящий сервер от подмены.
  \item Новая прошивка пишется в свободный слот, а рабочая остаётся нетронутой: при сбое обновления есть к чему вернуться.
\end{enumerate}

\anslesson{Урок 14. Диагностика}
\begin{enumerate}
  \item Адреса функций меняются при каждой сборке; чужой ELF даст правдоподобные, но ложные строки.
  \item О просадке питания. Искать нужно в источнике, кабеле и мощной нагрузке, а не в программе.
  \item Регулярно отдавать процессор: \code{vTaskDelay()}, ожидание очереди с тайм-аутом, никаких долгих циклов без пауз.
\end{enumerate}

\anslesson{Урок 15. Датчики}
\begin{enumerate}
  \item Питание, логические уровни, интерфейс, время и период опроса, предельные значения.
  \item Одно измерение занимает до двух секунд; при более частых запросах датчик возвращает старые данные или ошибку.
  \item ECHO выдаёт \SI{5}{V}, а вход платы рассчитан на \SI{3.3}{V}. Расстояние в сантиметрах: время эха в микросекундах × 0,0343 / 2.
\end{enumerate}

\anslesson{Урок 16. Силовая нагрузка}
\begin{enumerate}
  \item Мотору нужны сотни миллиампер и другое напряжение, а вывод выдерживает лишь десятки миллиампер при \SI{3.3}{V}.
  \item Он гасит выброс напряжения при отключении катушки. Диод ставят параллельно катушке, катодом к плюсу питания.
  \item Напряжение управляющего сигнала отсчитывается от земли; без общей земли драйвер «не понимает» уровни платы.
\end{enumerate}

\anslesson{Урок 17. Адресные светодиоды}
\begin{enumerate}
  \item Длительностью HIGH: короткий импульс --- ноль, длинный --- единица; длина бита одинакова.
  \item До $30 \cdot 60$ мА $= 1{,}8$ А. Нужен отдельный источник \SI{5}{V} и ограничение яркости в программе.
  \item Резистор гасит выбросы и отражения в линии данных, конденсатор сглаживает бросок тока при включении.
\end{enumerate}

\anslesson{Урок 18. HTTP, JSON, время, файлы}
\begin{enumerate}
  \item Плата узнаёт IP-адрес сервера через DNS, устанавливает соединение, отправляет запрос и ждёт код состояния и тело ответа.
  \item Ответ может прийти неполным или не в формате JSON; без проверки программа будет работать с мусором.
  \item NVS --- небольшие настройки «ключ--значение»; LittleFS --- файлы во Flash платы; SD-карта --- большие объёмы на съёмном носителе.
\end{enumerate}

\anslesson{Урок 19. Нативный USB}
\begin{enumerate}
  \item «UART» даёт только последовательный порт через микросхему-мост; нативный USB идёт прямо к чипу, который может стать клавиатурой, мышью и другими устройствами.
  \item \textsf{Tools → USB Mode: USB-OTG (TinyUSB)}.
  \item Отключить кабель от разъёма «USB», войти в загрузчик кнопками BOOT и RST и залить исправленную программу через «UART».
\end{enumerate}

\anslesson{Урок 20. Камера}
\begin{enumerate}
  \item Восемь линий данных, тактовый сигнал PCLK, сигналы кадра VSYNC и строки HREF, опорная частота XCLK. По SCCB (это I\textsuperscript{2}C) плата записывает настройки сенсора.
  \item Он возвращает буфер кадра драйверу. Без него свободные буферы закончатся, и следующий кадр получить не удастся.
  \item PSRAM (кадр хранится в ней), схема разделов Huge APP (программа с камерой большая) и USB CDC On Boot (чтобы видеть сообщения).
\end{enumerate}

\anslesson{Урок 21. Таймеры и конечные автоматы}
\begin{enumerate}
  \item Таймер срабатывает по прерыванию со строгим периодом, даже если основная программа занята; проверка \code{millis()} зависит от скорости \code{loop()}.
  \item Состояние --- что устройство делает сейчас (горит красный). Событие --- что произошло (истекло время, нажата кнопка). Переход --- смена состояния по событию (красный → красный с жёлтым).
  \item Вместе с состоянием нужно обновить время входа и выходы; если это сделано в одном месте, ничего нельзя забыть.
\end{enumerate}

\anslesson{Урок 22. Биты, байты и регистры}
\begin{enumerate}
  \item $32 + 8 + 4 + 1 = 45$; по четвёркам: 0010 = 2, 1101 = D, то есть \code{0x2D}.
  \item \code{x \&= \textasciitilde(1 << 5);}
  \item Пока биты вдвигаются, они «бегут» по выходам. Защёлка обновляет все восемь выходов одновременно, когда передача закончена.
\end{enumerate}

\anslesson{Урок 23. Органы управления}
\begin{enumerate}
  \item В момент спада сигнала A читают уровень B: HIGH --- одно направление, LOW --- другое.
  \item Кнопки стоят на пересечениях 4 строк и 4 столбцов; программа опрашивает строки по очереди.
  \item «Спокойный» уровень зависит от длины провода, влажности и питания, поэтому порог нельзя задать заранее.
\end{enumerate}

\anslesson{Урок 24. Звук}
\begin{enumerate}
  \item Пассивной нужен сигнал нужной частоты, активная пищит сама от постоянного напряжения. Если подать \SI{3.3}{V}: пассивная щёлкнет, активная запищит.
  \item Ровно в два раза.
  \item Пока идёт \code{delay()}, программа не реагирует на кнопки и датчики.
\end{enumerate}

\anslesson{Урок 25. Шаговый двигатель}
\begin{enumerate}
  \item Обычный мотор просто вращается и не знает положения; сервопривод держит заданный угол в пределах 180°; шаговый вращается сколько угодно, но шагами, которые можно считать.
  \item $4096 \cdot 45 / 360 = 512$ полушагов.
  \item Ротор не успевает за полем и пропускает шаги; программа продолжает считать и теряет настоящее положение вала.
\end{enumerate}

\anslesson{Урок 26. Акселерометр и гироскоп}
\begin{enumerate}
  \item \code{(int16\_t)((старший << 8) | младший)}.
  \item Малая ошибка скорости при накоплении складывается, и угол постепенно уплывает (дрейф).
  \item Доля акселерометра вырастет: дрейф исчезнет быстрее, но в угол попадут шум и вибрация.
\end{enumerate}

\anslesson{Урок 27. Автономное питание}
\begin{enumerate}
  \item Перезаряд, глубокий разряд и короткое замыкание повреждают аккумулятор и могут привести к возгоранию; плата защиты отключает его в этих случаях.
  \item Напряжение до \SI{4.2}{V} выше предела АЦП. Через большие резисторы течёт ничтожный ток, а малые сами разряжали бы аккумулятор.
  \item Около $1000 / 20 = 50$ часов, на практике на 20--30 \% меньше.
\end{enumerate}

\anslesson{Урок 28. Порядок в коде}
\begin{enumerate}
  \item Данные и функции собраны вместе и защищены от случайной порчи; для второй кнопки достаточно создать второй объект.
  \item В \code{.h} --- объявление класса (что он умеет), в \code{.cpp} --- тела функций (как он это делает).
  \item После каждого рабочего шага и перед рискованной правкой; в описании --- что теперь умеет программа.
\end{enumerate}

\anslesson{Урок 29. Цифровой звук: I²S и микрофон}
\begin{enumerate}
\item I²S передаёт аудио с тактами BCLK/WS; I²C --- адресованные команды по SDA/SCL.
\item Полезные данные и длина передаваемого слота различаются; важны знак и выравнивание.
\item Получены относительные единицы отсчётов. Для дБ SPL нужны чувствительность и акустическая калибровка.
\end{enumerate}

\anslesson{Урок 30. Точные импульсы: периферия RMT}
\begin{enumerate}
\item 250 мкс.
\item PCNT считает фронты, а не все длительности и порядок уровней.
\item До подтверждения окончания; раньше нельзя освобождать или изменять его.
\end{enumerate}

\anslesson{Урок 31. Счётчики и движение: PCNT и MCPWM}
\begin{enumerate}
\item 250 при точном окне 250 мс.
\item После обнуления на границе обычная разность показаний будет неверной.
\item Пауза между переключениями плеч моста, снижающая риск их одновременного открытия.
\end{enumerate}

\anslesson{Урок 32. Проводная сеть: TWAI/CAN и RS-485}
\begin{enumerate}
\item Нужен трансивер логических и дифференциальных уровней.
\item На двух физических концах, с учётом уже установленных резисторов.
\item Функция ставит кадр в очередь; факт приёма и обработки нужно подтвердить.
\end{enumerate}

\anslesson{Урок 33. Цветной экран: SPI, RGB, память и DMA}
\begin{enumerate}
\item 153600 байт; два полных буфера --- 307200 байт.
\item SPI записывает команды и пиксели в контроллер, RGB-панель требует потока с синхронизацией.
\item У драйверов разные требования к доступу, выравниванию и буферам.
\end{enumerate}

\anslesson{Урок 34. USB глубже: CDC, накопитель и роль Host}
\begin{enumerate}
\item Да: это поток байтов, границы задаёт приложение.
\item Файловые системы и кэши получают независимые несогласованные изменения.
\item Нужны подходящее питание VBUS, соединение, ограничение тока и стек класса.
\end{enumerate}

\anslesson{Урок 35. ESP-IDF и JTAG: увидеть выполнение}
\begin{enumerate}
\item Типы операндов определяют деление; дробь теряется до присваивания float.
\item Он содержит адреса, символы и связь со строками именно этой сборки.
\item Задачи останавливаются, а тайм-ауты удалённых устройств и часть периферии продолжаются.
\end{enumerate}

\anslesson{Урок 36. ULP: наблюдение, пока основные ядра спят}
\begin{enumerate}
\item Нужны компилятор RISC-V, встраивание образа и настройка RTC-памяти.
\item Нет: момент нажатия имеет произвольную фазу относительно опроса.
\item Ток моста USB, стабилизатора, LED и модулей может доминировать.
\end{enumerate}

\anslesson{Урок 37. Распознавание на устройстве: данные, модель и ESP-DL}
\begin{enumerate}
\item Обучение настраивает параметры, инференс применяет их.
\item Отсчёты зависимы: получается утечка информации и завышенная оценка.
\item Нужны ещё входы, промежуточные тензоры, кадры, стек и сеть.
\end{enumerate}

\anslesson{Урок 38. От макетки к устройству: запуск, восстановление и защита}
\begin{enumerate}
\item Нужны конфигурация загрузчика и самопроверка/подтверждение новой прошивки.
\item Первый проверяет подлинность запускаемого кода, вторая защищает содержимое внешней Flash.
\item Переключение при повторе меняет результат; повтор установки того же значения идемпотентен.
\end{enumerate}

\anslesson{Урок 39. ESPUI: веб-панель из C++}
\begin{enumerate}
  \item ID обозначает элемент интерфейса. GPIO обозначает электрический вывод; связь между ними задаёт программа.
  \item Callback может работать в сетевой задаче: он проверяет ввод и передаёт команду, а изменение модели и долгая работа выполняются в назначенном месте.
  \item Страница загружается по HTTP, команды идут по отдельному WebSocket. В ESPUI 2.2.4 его аутентификация зависит от отдельного флага, выключенного по умолчанию.
\end{enumerate}

\anslesson{Урок 40. WiFiManager: настройка сети}
\begin{enumerate}
  \item Первый позволяет войти во временную сеть настройки; второй нужен устройству для подключения к выбранному роутеру.
  \item В неблокирующем режиме процесс обслуживания продолжается после setup(), поэтому объект должен сохраняться.
  \item Два сервера не могут независимо занять один адрес и порт. Нужен общий сервер либо раздельные режимы с освобождением ресурсов.
\end{enumerate}

\anslesson{Урок 41. Итоговый проект}
\begin{enumerate}
  \item Она интересна автору, решаема за 4--6 недель и опирается на уже освоенные уроки.
  \item Это измеримая проверка, по которой видно, что проект готов. Список функций говорит, что устройство должно делать, а критерий --- как убедиться, что оно это делает.
  \item Входы, выходы, микроконтроллер, связь и память. Питание рисуют отдельно, потому что от него зависит работа всех блоков и именно в нём чаще всего ошибаются.
\end{enumerate}

\endgroup
